\chapter{Fractures}
\index{Fractures}
\label{sec:Fractures}

\section{Overview}
A fracture is a break in the continuity of bone. Common fracture types include distal radius (Colles' fracture), hip (intertrochanteric or subcapital femoral neck), vertebral compression, scaphoid, and clavicle fractures. This condition results from acute injury or exposure to harmful agents.
\section{Causes}
This condition happens because of mechanical force exceeding bone strength.     
\section{Risk Factors}

\begin{itemize}[noitemsep]
  \item Genetic predisposition and family history
  \item Advancing age
  \item Lifestyle factors (diet, exercise, smoking, alcohol)
  \item Environmental exposures
  \item Underlying medical conditions
  \item Medication use
\end{itemize}
\section{Signs and Symptoms}

\subsection{Common symptoms}
\begin{itemize}[noitemsep]
  \item You may experience pain, deformity, swelling, and loss of function. 
  \end{itemize}

\subsection{Emergency warning signs}
\begin{itemize}[noitemsep]
  \item Severe or worsening symptoms of fractures require immediate medical evaluation
\end{itemize}
\section{Progression and Stages}
Fractures are classified as open (compound, with skin disruption) or closed (simple), and by pattern (transverse, oblique, spiral, comminuted, greenstick, impacted, avulsion). The condition may progress through different stages of severity over time. Early stages often involve mild or subtle symptoms that may be overlooked. Moderate stages involve more noticeable symptoms affecting daily function. Advanced stages may involve significant impairment and complications.
\section{Complications}
\begin{itemize}[noitemsep]
  \item In most cases, the outlook is good
  \item Hip fractures in the elderly are orthopedic emergencies requiring surgical fixation within 24--48 hours to reduce mortality and complications.
  \item Disease progression leading to worsening symptoms
  \item Secondary infections or injuries
  \item Side effects of treatment
  \item Reduced quality of life
  \item Emotional and psychological impact
\end{itemize}
\section{Diagnosis}
\begin{itemize}[noitemsep]
  \item Doctors diagnose this using X-rays (at least two views)
  \item CT for complex fractures, MRI for occult fractures and ligamentous injury.
  \item Comprehensive medical history and physical examination
  \item Laboratory tests to assess organ function
  \item Imaging studies to visualize affected structures
  \item Specialized testing depending on the affected system
  \item Consultation with relevant specialists
\end{itemize}
\section{Prevention}
\begin{itemize}[noitemsep]
  \item Maintain a healthy lifestyle with balanced diet and exercise
  \item Avoid known risk factors
  \item Attend regular medical check-ups for early detection
  \item Follow recommended screening guidelines
\end{itemize}
\section{Treatment Options}
Medications to manage symptoms and address underlying mechanisms Lifestyle modifications and self-care measures Physical therapy and rehabilitation Surgical intervention when indicated Regular monitoring and follow-up care Patient education and support services

\begin{itemize}[noitemsep]
  \item Medications to manage symptoms and address underlying mechanisms
  \item Lifestyle modifications and self-care measures
  \item Physical therapy and rehabilitation
  \item Surgical intervention when indicated
  \item Regular monitoring and follow-up care
\end{itemize}
\section{Living with It}
Treatment aims for reduction (realigning fracture fragments), fixation (maintaining alignment during healing), and rehabilitation

\begin{itemize}[noitemsep]
  \item Closed reduction with casting is appropriate for stable, minimally displaced fractures
  \item Open reduction with internal fixation (ORIF) is indicated for displaced, unstable, or intra-articular fractures.
  \item Follow your treatment plan as prescribed
  \item Maintain regular follow-up with your healthcare provider
  \item Make necessary lifestyle adjustments
  \item Seek support from family, friends, or support groups
  \item Stay informed about your condition
\end{itemize}
\section{Outlook}
The outlook varies depending on the specific condition, severity at diagnosis, and response to treatment. Early intervention generally leads to better outcomes. Many conditions can be effectively managed with appropriate treatment. Regular follow-up care is important for monitoring and treatment adjustment.
